
\documentclass[11pt,a4paper]{article}
\usepackage[a4paper,margin=2.05cm]{geometry}
\usepackage[spanish,es-noquoting]{babel}
\usepackage[T1]{fontenc}
\usepackage[utf8]{inputenc}
\usepackage{lmodern,microtype}
\usepackage{amsmath,amssymb,amsthm,mathtools}
\usepackage{booktabs,longtable,array,enumitem}
\usepackage{xcolor,hyperref}
\pagecolor{white}\color{black}
\setlength{\parindent}{0pt}
\setlength{\parskip}{0.65em}
\title{\bfseries N72 --- Supervivencia coinductiva y novena puerta\\
\large Del selector local al inverso límite de frontera}
\author{HMT / auditoría técnica}
\date{}
\begin{document}
\maketitle

\begin{abstract}
N72 fortalece N71. El selector local de pares \((u,\delta)\) por firma de frontera no siempre es único. En lugar de cerrar en una sola frontera, se aplica supervivencia coinductiva: una rama es válida si admite extensiones sucesivas que satisfacen cilindro exacto y firma de frontera. Se audita \(t=5,\ldots,25\) hasta profundidad \(9\). Todas las ambigüedades se resuelven en profundidad menor o igual que \(3\), y la rama real sobrevive hasta la novena frontera futura. La novena puerta actúa como confirmación de estabilidad, no como aparición de nuevas ramas.
\end{abstract}

\section{Definición}

Sea \(S_t\) el conjunto local de ramas seleccionadas por:
\[
\text{cilindro exacto}+\text{firma de frontera}.
\]
Una rama \(s\in S_t\) sobrevive a profundidad \(h\) si existe una cadena:
\[
s=s_t,s_{t+1},\ldots,s_{t+h-1}
\]
tal que cada \(s_{j+1}\) pertenece al selector de frontera sobre el prefijo extendido por \(s_j\).

La rama coinductiva ideal es la que sobrevive para todo \(h\).

\section{Auditoría}

\[
\begin{array}{c|c|c|c|c|c|c}
t&\#S_t&h_{\rm res}&h1&h2&h3&h9\\
\hline
5 & 2 & 2 & 2 & 1 & 1 & 1 \\
6 & 6 & 3 & 6 & 3 & 1 & 1 \\
7 & 11 & 3 & 11 & 3 & 1 & 1 \\
8 & 3 & 3 & 3 & 2 & 1 & 1 \\
9 & 3 & 2 & 3 & 1 & 1 & 1 \\
10 & 11 & 2 & 11 & 1 & 1 & 1 \\
11 & 3 & 2 & 3 & 1 & 1 & 1 \\
12 & 5 & 2 & 5 & 1 & 1 & 1 \\
13 & 3 & 2 & 3 & 1 & 1 & 1 \\
14 & 3 & 2 & 3 & 1 & 1 & 1 \\
15 & 5 & 2 & 5 & 1 & 1 & 1 \\
16 & 11 & 2 & 11 & 1 & 1 & 1 \\
17 & 6 & 2 & 6 & 1 & 1 & 1 \\
18 & 1 & 0 & 1 & 1 & 1 & 1 \\
19 & 1 & 0 & 1 & 1 & 1 & 1 \\
20 & 5 & 3 & 5 & 5 & 1 & 1 \\
21 & 2 & 3 & 2 & 2 & 1 & 1 \\
22 & 1 & 0 & 1 & 1 & 1 & 1 \\
23 & 3 & 3 & 3 & 2 & 1 & 1 \\
24 & 3 & 2 & 3 & 1 & 1 & 1 \\
25 & 6 & 3 & 6 & 6 & 1 & 1 \\
\end{array}
\]

\section{Novena puerta}

En \(t=9\), que corresponde a \(K=9\), se tiene:
\[
\#S_9=3.
\]
A profundidad \(1\):
\[
3\text{ ramas sobreviven}.
\]
A profundidad \(2\):
\[
1\text{ rama sobrevive}.
\]
A profundidad \(9\):
\[
1\text{ rama sobrevive}.
\]
Por tanto:
\[
\boxed{
\text{la novena puerta confirma estabilidad; no introduce falsa rama nueva.}
}
\]

\section{Dictamen}

\[
\boxed{
\text{verde: todas las ambigüedades auditadas se resuelven por supervivencia finita.}
}
\]
\[
\boxed{
\text{verde: la profundidad máxima necesaria en }t=5,\ldots,25\text{ es }3.
}
\]
\[
\boxed{
\text{verde: la rama real sobrevive hasta la novena frontera futura.}
}
\]
\[
\boxed{
\text{amarillo: esto es auditoría finita, no demostración infinita completa.}
}
\]

\section{Lectura}

La supervivencia no es un añadido externo: es el modo natural de leer números HMT como objetos de límite inverso. El número no es una cinta puntual; es la rama que sobrevive bajo todas las fronteras compatibles.
\end{document}
